% Page budget: 2.90 pages | Owns: all quantitative evidence, figure and table claims, criteria, verdicts, and comparison to the black box.
% WRITING GUIDE
% Purpose: Answer the research aspects with frozen quantitative evidence, including negative and conditional results.
% Start here: Final run artefacts and generated figure data; use docs/gantry-augmentation-problem-log.md only to locate and interpret them.
% Supporting material: Quinten's 18 September outline feedback, final configuration manifests, checkpoint reports, controller checks, OBC monitoring, ablations, and black-box outputs.
% Research-plan anchor: Evaluate each promised aspect against its original criterion; name any missing or changed experiment as a departure.
% Current decisions: Every result states expectation, criterion, verdict, and mechanism under matched data and optimisation budgets.
% Choices to justify: Comparator capacity, matched budget, checkpoint rule, one primary metric, representative ASMPT motion, settling interval, controller-transfer case, figure aggregation, and treatment of neural seeds separately from phase realizations.
% Implementation audit: Trace figures and tables to gantry_dynamic/evaluation.py, diagnostics.py, obc_diagnostics.py, lpv_lfr_baseline evaluation scripts, final configs, and saved artefacts.
% Reference check: Results normally cite methods or benchmark definitions, not external authority for the measured outcome; keep internal provenance in captions, configs, and repository records.
% Cautions: Do not promote preregistrations, implementation plans, diagnostic-only runs, or superseded numbers to final evidence.
% Planned figures: Use only the result figures and tables named in the subsection bullets; combine panels where they support one claim.
% Section is finished when: Every number is reproducible, every comparison is fair, and every claimed mechanism has an intervention or derivation behind it.
\section{Results}

\subsection{Baseline and closed-loop formulation}
\begin{itemize}
\item Expectation and criterion: the LPV-LFR and original physics equations agree numerically across the payload range before model-error comparisons are meaningful
\item Verdict and reason: report inertia invertibility, LFR/RK4 equivalence, and the remaining baseline error caused by the truth-only absorber
\item Figure: baseline and oracle output errors versus time for representative payload positions; supports the claim that the omitted mode is systematic and position dependent
\item Table: baseline free-run physical RMS for LPV and frozen-LTI models across held-out $Y$; tests the research plan's Aspect 1 rather than assuming LPV superiority
\item Expectation and criterion: direct and residual closed-loop forms match at numerical precision when controller assumptions hold
\item Figure: recorded input, controller correction, and model error over one rollout; supports the interpretation that feedback can suppress visible output error while the open-loop model remains wrong
\end{itemize}
\todo{The LPV-versus-frozen-LTI comparison is required by the plan but may not yet have a final matched result; run it or narrow Aspect 1 explicitly.}

\subsection{Predictive benefit and additional-state contribution}
\begin{itemize}
\item Expectation and criterion: augmentation must reduce held-out long-horizon error, not only the short training loss, relative to the same encoder-initialised baseline
\item Verdict and reason: report the best validated improvement and the gap between window fit and free-run accuracy caused by integrated velocity bias
\item Figure: validation RMS against optimiser updates for baseline and augmented arms, with window and 12 s metrics; supports or rejects objective alignment
\item Figure: error spectra before and after augmentation with absorber pole and anti-resonance marked; tests whether the learned correction targets the missing band
\item Figure: clamp each additional state and the full added-state vector to zero during rollout; quantifies what the trained latent dynamics add or harm
\item Table: standard model, no-added-state ablation, zero-initialised latent path, and longer encoder history; connects structure changes to predictive benefit
\item State the observed weak or harmful added-state contribution as a result if the final run does not reverse it, rather than inferring usefulness from state magnitude
\end{itemize}
\todo{Are state Gramians useful for explaining each added state's contribution, or would intervention-based ablation be more direct and less linearisation-dependent?}

\subsection{Parameter interpretability and orthogonality}
\begin{itemize}
\item Expectation and criterion: OBC should reduce physical combination error without materially worsening held-out predictive fit relative to unconstrained joint estimation
\item Verdict and reason: the completed arms reach 7.20\% no projection, 7.69\% tangent OBC, and 8.47\% unfinished affine OBC from a common 9.49\% start, so aggregate recovery does not establish OBC superiority
\item Figure: validation fit versus combination error over training for all three arms; supports the claim that near-equal fit permits materially different physical parameters
\item Figure: per-parameter error at matched budget; shows which combinations projection protects and which near-flat directions redistribute error
\item Table: basis rank, condition, overlap, final and minimum parameter error, and validation RMS; links algorithm health to the recovery verdict
\item Figure: $\rho^*$ and predicted versus observed parameter bias by singular direction; supports the conditional result that true recovery is limited by non-orthogonal missing physics
\item Explain why the 4.96\% predicted floor accounts for part of the residual squared error but does not match the observed per-parameter pattern or prove convergence
\end{itemize}
\todo{Should the incomplete affine arm be retained, rerun to a matched budget, or reported only as exploratory evidence?}

\subsection{Generalisation and black-box comparison}
\begin{itemize}
\item Expectation and criterion: improvement must progress from held-out position/phase interpolation to ASMPT-relevant motion, controller transfer, and quantified operating-range extrapolation without retuning
\item Verdict and reason: separate completed simulation evidence from pending Telica results and explain any failure by excitation coverage, initialisation, or objective mismatch
\item Figure: tracking overview plus a dedicated 200--400 ms settling zoom for baseline, augmented model, and black box on an approved operational trajectory
\item Table: primary RMS, fixed-window settling RMS, parameter availability, and training data for baseline, augmented, and black-box models on every completed validation rung
\item Figure or table: train with one controller and evaluate the frozen models with a credible changed controller; pair this with a plant-domain check for feedback masking
\item Figure: error versus multisine amplitude, stroke, acceleration, scheduling position, or state-input distance from the training set; quantify extrapolation rather than calling all held-out data extrapolation
\item Report the black box's early improvement and later free-run degradation as a negative comparator result if no stable matched run supersedes it
\item Aggregate the final learned-model comparison across neural initialization seeds and show a representative, not best, time trace; report phase realizations separately
\end{itemize}
\todo{Freeze which validation rungs are complete enough for a final claim; omit rather than speculate about unfinished controller, noise, or hardware experiments.}
